% !TeX root = main.tex
\chapter{事件与概率}

\subsection{Secret Santa}

假设你在一个有56名学生的班级里。老师要求每名学生抽取一名学生的名字，秘密地给另一名学生准备礼物。

假设很不幸\UseVerb{sad}，你时运不济、孤寡一生，抽到了自己的名字。除了自怨自哀之外，你不免好奇这样非酋的事件发生的概率。
班里互赠礼物可以看作是班内元素的一个排列，计算整个班内出至少这么一个倒霉学生的概率，也就是经典的错排问题。不难计算出，
抽到自己名字的概率是\(\frac{1}{56}\)。
其概率约为\(\frac{D_{1}(56)}{56!}\approx \e^{-1} =36.8\%\).
% todo: 错排数的计算，容斥原理

而有些人就是天生的欧皇体质，两个人相互抽到了对方的名字，在天愿作比翼鸟，在地愿为连理枝。班里出这样一对幸运儿的概率又是多大呢？

对于每对学生来说，他们相互抽到对方名字的概率是\(\frac{1}{n(n-1)}\)，
在此例中为\(\frac{1}{56 \times 55}\)，约是万分之三的概率。
而班里至少有一对幸运儿的概率的计算就比较麻烦了\cite[p.~24]{arndt2010generating}。

考虑具有\(n\) 名学生的班级。由\(n\)的有限性，以任意学生为起点，不断的找到送礼物的对象，总能形成一个轮换。考虑排列的轮换分解，
记最小轮换大小为\(m\) ，则班内不存在一对幸运儿，等价于\(m\geq 3\)。
记\(D_{m}(n)\) 为所有排列中最小轮换大小大于等于\(m\) 的排列数。

对于一个满足条件的递推，考虑其第\(n\) 位元素。
\begin{itemize}
    \item 若第\(n\) 位元素所在轮换大小大于\(m\) ，则去掉该元素，剩余排列仍满足条件，
        其数目为\(D_{m}(n-1)\)。第\(n\) 位元素可作为此\(n-1\) 元素中任意元素的后继。
        总计\((n-1)D_{m}(n-1)\)。
    \item 若第\(n\) 位元素所在轮换大小等于\(m\) ，则去掉该轮换，剩余排列仍满足条件，
        其数目为\(D_{m}(n-m)\)。
        在\(n\) 个元素中选出\(m\) 个元素组成轮换的方式有\(\combination{n-1}{m-1}\) 种，
        其全排列的方式有\((m-1)! \) 种。总计\(\combination{n-1}{m-1}(m-1)!D_{m}(n-m)\)。
\end{itemize}
因此有递推关系式：
\begin{align*}
    D_{m}(n)&=(n-1)D_{m}(n-1)+\combination{n-1}{m-1}(m-1)!D_{m}(n-m) \\
    &= (n-1)D_{m}(n-1)+\permutation{n-1}{m-1}D_{m}(n-m)
\end{align*}
前若干项为\footnote{可见\href{https://oeis.org/A038205}{OEISA038205}}：
\begin{lstlisting}[language=Mathematica]
d[m_, n_] := d[m, n] = (n - 1) d[m, n - 1] +
    FactorialPower[n - 1, m - 1] d[m, n - m]
d[m_, 0] := 1
d[m_, k_ /; k < 0] := 0
d[3, 56]/56! // N
\end{lstlisting}
\begin{table}[H]
    \centering
    \begin{tabular}{|*{12}{c|}}
        \hline
        \(n\)        & 0 & 1 & 2 & 3 & 4 & 5  & 6   & 7    & 8    & 9     & 10 \\
        \hline
        \(D_{3}(n)\) & 1 & 0 & 0 & 2 & 6 & 24 & 160 & 1140 & 8988 & 80864 & 809856\\
        \hline
    \end{tabular}
\end{table}
故产生幸运儿的概率为：

\[
    1-\frac{D_{3}(56)}{56!} \approx 1 - \e^{-\frac{3}{2}}   = 77.7\%
\]
渐进分析表明\cite[p.~176]{Generatingfunctionology}:
\[
    \lim_{n \to \infty} D_{m}(n)/n! = \exp(-H_{m-1})
\]
% todo: calculate the limit using generating functions

班内产生爱情的概率高达77.7\%，但每对学生却只有万分之三的概率能遇到另一半。这又何尝不像极了我们的人生呢？\UseVerb{sad}
